\documentclass[11pt,a4paper]{article}

\usepackage[utf8]{inputenc}
\usepackage[T1]{fontenc}
\usepackage[spanish]{babel}
\usepackage{geometry}
\usepackage{booktabs}
\usepackage{enumitem}
\usepackage{listings}
\usepackage{xcolor}
\usepackage{hyperref}
\usepackage{parskip}
\usepackage{titlesec}
\usepackage{fancyhdr}
\usepackage{tikz}
\usetikzlibrary{arrows.meta,positioning}

\geometry{margin=2.5cm}

\definecolor{codegray}{rgb}{0.95,0.95,0.95}
\definecolor{codegreen}{rgb}{0,0.5,0}
\definecolor{codepurple}{rgb}{0.4,0,0.6}
\definecolor{uvgblue}{rgb}{0.1,0.22,0.36}

\lstdefinestyle{code}{
  backgroundcolor=\color{codegray},
  basicstyle=\ttfamily\small,
  keywordstyle=\color{codepurple}\bfseries,
  commentstyle=\color{codegreen},
  breaklines=true,
  frame=single,
  framerule=0pt,
  xleftmargin=6pt,
  xrightmargin=6pt,
  aboveskip=6pt,
  belowskip=6pt,
}
\lstset{style=code}

\pagestyle{fancy}
\fancyhf{}
\rhead{\small CC3071 \textperiodcentered{} UVG 2026-I}
\lhead{\small Proyecto 02: YAPar}
\cfoot{\thepage}

\titleformat{\section}{\large\bfseries\color{uvgblue}}{}{0em}{}[\titlerule]
\titleformat{\subsection}{\normalsize\bfseries}{}{0em}{}

% -----------------------------------------------------------------------
\begin{document}

\begin{center}
  {\Large\bfseries Proyecto 02: YAPar}\\[4pt]
  {\large Generador de Analizadores Sintácticos SLR}\\[8pt]
  \small Universidad del Valle de Guatemala\\
  Facultad de Ingeniería: CC3071 Diseño de Lenguajes de Programación\\
  Semestre I, 2026\\[8pt]
  \begin{tabular}{ll}
    Diego Lopez     & \#23747 \\
    Jennifer Toxcon & \#21276 \\
    Roberto Barreda & \#23354 \\
  \end{tabular}\\[6pt]
  \small Repositorio: \url{https://github.com/bar23354/PRY2-DLP}
\end{center}

\hrule
\vspace{10pt}

% -----------------------------------------------------------------------
\section{Descripción del proyecto}

YAPar es un generador de analizadores sintácticos SLR implementado en Python.
Toma como entrada una gramática libre de contexto en lenguaje YAPar
(\texttt{.yalp}) y una especificación léxica YALex (\texttt{.yal}), valida la
consistencia de tokens entre ambos archivos y produce:

\begin{itemize}[noitemsep]
  \item Un archivo \texttt{theParser.py} completamente autónomo que implementa
        el analizador sintáctico generado.
  \item El autómata LR(0) visualizado como imagen (PNG o PDF).
  \item Detección y reporte de conflictos SLR, errores léxicos y errores
        sintácticos con número de línea y columna.
\end{itemize}

% -----------------------------------------------------------------------
\section{Integrantes}

\begin{tabular}{lll}
  \toprule
  \textbf{Nombre} & \textbf{Carné} & \textbf{Responsabilidad principal} \\
  \midrule
  Diego Lopez     & 23747 & Lector \texttt{.yalp}, modelos de gramática, tests \\
  Jennifer Toxcon & 21276 & Generador SLR, FIRST/FOLLOW, autómata LR(0) \\
  Roberto Barreda & 23354 & CLI, runtime del parser, generación de código, GUI \\
  \bottomrule
\end{tabular}

% -----------------------------------------------------------------------
\section{Requisitos}

\begin{itemize}[noitemsep]
  \item Python 3.10 o superior.
  \item Dependencias: \lstinline|pip install -r requirements.txt|
        (pytest, customtkinter, Pillow, graphviz, matplotlib, networkx).
  \item Graphviz opcional para mejor resolución del diagrama LR(0):
        \lstinline|winget install Graphviz.Graphviz|.
\end{itemize}

% -----------------------------------------------------------------------
\section{Diseño del software}

\subsection{Arquitectura por capas}

El proyecto sigue una arquitectura en capas separadas por responsabilidad.
Cada módulo expone una interfaz pública mínima y no conoce los detalles
de implementación de los demás.

\begin{center}
\begin{tabular}{lp{9cm}}
  \toprule
  \textbf{Módulo} & \textbf{Responsabilidad} \\
  \midrule
  \texttt{src/grammarModel.py}  &
    Define los tipos inmutables \texttt{Production} y \texttt{Grammar} usados
    como contrato entre todas las capas. \\
  \texttt{src/yalpReader.py}    &
    Elimina comentarios \texttt{/* */}, separa secciones con \texttt{\%\%},
    procesa \texttt{\%token} e \texttt{IGNORE}, parsea producciones y valida
    nombres y símbolos. Sin expresiones regulares. \\
  \texttt{src/yalexReader.py}   &
    Lee un archivo \texttt{.yal} y extrae los nombres de tokens definidos
    después de \texttt{->}. Valida que los tokens del \texttt{.yalp} existan
    también en el lexer. \\
  \texttt{src/slrGenerator.py}  &
    Implementa el algoritmo SLR completo: aumentación, FIRST, FOLLOW,
    clausura, goto, construcción de estados LR(0), tablas ACTION/GOTO y
    detección de conflictos. \\
  \texttt{src/lr0Visualizer.py} &
    Genera el diagrama del autómata LR(0) usando Graphviz si está disponible
    en el PATH; si no, usa matplotlib + networkx como respaldo. \\
  \texttt{src/parserRuntime.py} &
    Implementa el ciclo principal del parser: shift, reduce, accept y reporte
    de errores léxicos y sintácticos con posición. \\
  \texttt{src/codeGenerator.py} &
    Serializa las tablas ACTION/GOTO en texto Python y las inserta en una
    plantilla que produce \texttt{theParser.py} autocontenido. \\
  \texttt{src/lexerAdapter.py}  &
    Tokenizador simple que mapea palabras del archivo de entrada a tipos de
    token; produce tokens en formato dict compatible con el runtime. \\
  \texttt{src/guiApp.py}        &
    Interfaz gráfica con siete pestañas: Gramática, FIRST/FOLLOW, Estados
    LR(0), Tablas, Autómata, Generar Parser y Probar Parser. \\
  \texttt{yapar.py}             &
    Punto de entrada CLI: valida argumentos, orquesta el pipeline completo
    y reporta errores al usuario. \\
  \bottomrule
\end{tabular}
\end{center}

\subsection{Flujo de datos}

El pipeline principal sigue este orden:

\begin{enumerate}[noitemsep]
  \item \textbf{Lectura}: \texttt{parseYalp()} lee el \texttt{.yalp} y devuelve
        un objeto \texttt{Grammar} con tokens, tokens ignorados, producciones y
        símbolo inicial.
  \item \textbf{Validación léxica}: \texttt{validateTokensInYalex()} verifica que
        cada token del \texttt{.yalp} esté definido en el \texttt{.yal}.
  \item \textbf{Generación SLR}: \texttt{buildSlrParserTable()} aumenta la
        gramática, calcula FIRST/FOLLOW, construye los estados LR(0) y llena
        las tablas ACTION y GOTO. Si hay conflictos lanza \texttt{SLRConflictError}.
  \item \textbf{Visualización}: \texttt{generateLr0Diagram()} produce el PNG del
        autómata.
  \item \textbf{Generación de código}: \texttt{generateParserFile()} escribe
        \texttt{theParser.py} con las tablas serializadas.
  \item \textbf{Ejecución}: \texttt{theParser.py} tokeniza el archivo de entrada
        con \texttt{tokenize()} y ejecuta el ciclo shift/reduce/accept.
\end{enumerate}

\subsection{Algoritmo SLR: implementación detallada}

\paragraph{Aumentación.} Se añade la producción $S' \to S$ como producción~0.
El símbolo aumentado $S'$ actúa como nuevo símbolo inicial y permite que el
estado de aceptación tenga un único ítem kernel.

\paragraph{Conjuntos FIRST y FOLLOW.} Se calculan iterativamente hasta punto
fijo. \texttt{computeFirst} soporta producciones vacías
($\varepsilon$); \texttt{computeFollow} añade \texttt{\$} a
$\text{FOLLOW}(S')$ y propaga siguiendo las reglas estándar.

\paragraph{Clausura y goto.} \texttt{closure} expande un conjunto de ítems
LR(0) añadiendo todas las producciones de los no-terminales que aparecen
después del punto. \texttt{goto} mueve el punto sobre un símbolo y aplica
clausura al resultado.

\paragraph{Construcción de estados.} \texttt{buildLr0States} parte del estado
$I_0 = \text{closure}(\{S' \rightarrow \cdot S\})$ y aplica \texttt{goto} con todos
los símbolos de la gramática para obtener nuevos estados, evitando duplicados
mediante un mapa de conjuntos de ítems a identificadores de estado.

\paragraph{Tablas ACTION y GOTO.}
\begin{itemize}[noitemsep]
  \item Si $A \to \alpha \cdot a \beta$ y $a$ es terminal: $\text{ACTION}[s,a] = \text{shift}(s')$.
  \item Si $A \to \alpha \cdot$ y $A \neq S'$: $\text{ACTION}[s,t] = \text{reduce}(i)$
        para todo $t \in \text{FOLLOW}(A)$.
  \item Si $S' \to S \cdot$: $\text{ACTION}[s,\$] = \text{accept}$.
  \item $\text{GOTO}[s, A] = s'$ para cada no-terminal $A$.
\end{itemize}
Un conflicto ocurre cuando se intenta escribir dos acciones distintas en la
misma celda; en ese caso \texttt{SLRConflictError} detiene el proceso.

\subsection{Parser generado: theParser.py}

\texttt{codeGenerator.py} serializa las tablas como literales Python y las
inserta en una plantilla de texto. El archivo resultante:

\begin{itemize}[noitemsep]
  \item Contiene \texttt{TOKENS}, \texttt{IGNORE\_TOKENS}, \texttt{PRODUCTIONS},
        \texttt{ACTION} y \texttt{GOTO} como diccionarios y listas Python.
  \item Define \texttt{\_splitWithColumns()} para tokenizar sin expresiones
        regulares: recorre la línea carácter a carácter y registra la columna
        de inicio de cada palabra.
  \item Define \texttt{parse(tokens)}: implementa el ciclo
        shift/reduce/accept con una pila de estados enteros.
  \item Se ejecuta directamente con \lstinline|python theParser.py input.txt|.
\end{itemize}

\subsection{Documentación del código}

Todos los módulos siguen las mismas convenciones:

\begin{itemize}[noitemsep]
  \item \textbf{Nombres}: \texttt{camelCase} para funciones y variables;
        \texttt{PascalCase} para clases y dataclasses.
  \item \textbf{Dataclasses}: \texttt{Production} e \texttt{Item} son
        \texttt{frozen=True} para poder usarlos en conjuntos y como claves
        de diccionario.
  \item \textbf{Sin efectos secundarios}: las funciones del pipeline SLR son
        puras; toda la mutabilidad queda en \texttt{guiApp.py}.
  \item \textbf{Sin expresiones regulares}: ningún módulo importa \texttt{re}.
        El reconocimiento de patrones se implementa con operaciones de strings
        (\texttt{find}, \texttt{split}, \texttt{startswith}) y recorridos de
        caracteres.
\end{itemize}

% -----------------------------------------------------------------------
\section{Cómo ejecutarlo}

\subsection{CLI}

\begin{lstlisting}[language=bash]
# Generar el parser
python yapar.py examples/medium/parser.yalp \
       -l examples/medium/lexer.yal \
       -o theParser.py

# Ejecutar el parser generado
python theParser.py examples/medium/input_valid.txt
python theParser.py examples/medium/input_error.txt
\end{lstlisting}

\subsection{Interfaz gráfica}

\begin{lstlisting}[language=bash]
python src/guiApp.py
\end{lstlisting}

La GUI ofrece las mismas capacidades que el CLI desde una sola ventana:
carga de archivos, validación YAPar/YALex, análisis de gramática,
visualización del autómata LR(0), generación de \texttt{theParser.py} y
prueba del parser con un archivo de entrada.

% -----------------------------------------------------------------------
\section{Archivos de prueba}

Se proveen tres grupos de archivos, uno por nivel de complejidad.
Cada grupo contiene: especificación de lexer (\texttt{.yal}),
especificación de parser (\texttt{.yalp}), entrada válida e
entrada con error.

\begin{center}
\begin{tabular}{lll}
  \toprule
  \textbf{Complejidad} & \textbf{Gramática} & \textbf{Directorio} \\
  \midrule
  Baja   & Suma de números
          (\texttt{expr $\to$ NUM $\mid$ expr PLUS NUM})
        & \texttt{examples/low/} \\
  Media  & Aritmética con precedencia ($+$, $\times$, paréntesis)
        & \texttt{examples/medium/} \\
  Alta   & Asignaciones e instrucción \texttt{IF\ldots THEN\ldots END}
        & \texttt{examples/high/} \\
  \bottomrule
\end{tabular}
\end{center}

\end{document}
